\subsubsection{实验一 AI交互记录}\label{ux5b9eux9a8cux4e00-aiux4ea4ux4e92ux8bb0ux5f55}

\begin{enumerate}
\def\labelenumi{\arabic{enumi}.}
\tightlist
\item
  需求拆解：读取《金融数据挖掘实验指导书》，确认实验一目标为金融新闻文本清洗、情感标注、按5个交易日生成周度情绪指标，并输出CSV与数据库。
\item
  数据归档：将老师发布的实验指导书和配套数据归档到 \texttt{data/experiment1/original/}，将新闻数据与沪深300价格数据放入 \texttt{data/experiment1/raw/}。
\item
  数据清洗：编写Python代码读取新闻表，筛选2014-10至2015-10样本，去除空文本、重复文本和非交易日新闻。
\item
  情感标注：优先使用本地 \texttt{yiyanghkust/finbert-tone-chinese} 金融中文BERT模型进行正面/中性/负面三分类；若模型文件缺失，则使用透明金融词典规则作为可运行兜底。
\item
  周度指标：按沪深300交易日历每5个交易日形成一周，统计 \texttt{WeekPositive}、\texttt{WeekNeutral}、\texttt{WeekNegative}、\texttt{NewsCount} 和 \texttt{P\_t}。
\item
  数据库存储：将周度情绪表、标注样本、羊群指标和建模数据写入SQLite数据库，并同步导出CSV。
\item
  质量检查：检查样本日期、去重数量、交易日过滤结果、周度字段完整性和数据库输出。
\item
  报告生成：输出Markdown报告、LaTeX/PDF报告、核心图表、数据库文件和代码附件，用于提交。
\end{enumerate}
